% ===========================================================================
% Section: Gaussian Reconstruction
% ===========================================================================
\section{Gaussian Splatting Reconstruction}
\label{sec:gaussian_reconstruction}

The reconstruction pipeline preserves the core algorithm of
\gaussianlic~\cite{lang2026_gaussianlic2} while replacing the ROS-based
data ingestion with the offline sensor source layer. This section
describes the four stages: initialization, frame addition, map extension,
and map optimization.

\subsection{Gaussian Initialization}
\label{subsec:init}

The Gaussian map is initialized from the first keyframe's fused point
cloud (Figure~\ref{fig:gaussian_init}).

% ===========================================================================
% Figure: Gaussian initialization from fused point cloud
% Fused point cloud → SH colour → scale estimation → 3DGS parameters
% ===========================================================================
\begin{figure*}[t]
\centering
\begin{tikzpicture}[node distance=0.8cm and 1.0cm]

% ---- Input: Fused point cloud ----
\node[blockgreen, minimum width=3.5cm, minimum height=1.2cm]
  (cloud) {Fused Point Cloud\\$\PointCloud = \{(\mathbf{p}_i, \mathbf{c}_i, d_i)\}$};

% ---- Branch 1: SH colour ----
\node[block, below left=1.2cm and 0.3cm of cloud, minimum width=3.2cm, fill=blue!10]
  (sh) {Spherical Harmonics\\$\mathbf{f}_i^{\text{dc}} = \mathrm{RGB2SH}(\mathbf{c}_i)$};

% ---- Branch 2: Scale estimation ----
\node[block, below right=1.2cm and 0.3cm of cloud, minimum width=3.2cm, fill=orange!10]
  (scale) {Scale Estimation\\$s_i = \log\!\left(\alpha \cdot d_i / f\right)$};

% ---- Branch 3: Rotation & opacity ----
\node[block, below=3.0cm of cloud, minimum width=3.2cm, fill=purple!10]
  (rot) {Rotation \& Opacity\\$\mathbf{q}_i = [1,0,0,0]$, $o_i = \sigma^{-1}(0.1)$};

% ---- Merge: 3DGS parameters ----
\node[blockgreen, below=2.0cm of rot, minimum width=5cm, minimum height=1.2cm]
  (params) {3D Gaussian Parameters\\
            $\{\mathbf{x}_i, \mathbf{f}_i, s_i, \mathbf{q}_i, o_i\}_{i=1}^{N}$};

% ---- Output: Gaussian map ----
\node[block, below=of params, minimum width=5cm, fill=gray!10]
  (map) {Gaussian Map $\GaussianMap$};

% ---- Arrows ----
\draw[arrow] (cloud) -- (sh);
\draw[arrow] (cloud) -- (scale);
\draw[arrow] (cloud) -- (rot);
\draw[arrow] (sh)    -- (params.north -| sh.south);
\draw[arrow] (scale) -- (params.north -| scale.south);
\draw[arrow] (rot)   -- (params);
\draw[arrow] (params) -- (map);

% ---- Skybox (optional, side branch) ----
\node[blockred, right=2.0cm of cloud, minimum width=2.8cm, minimum height=1.0cm]
  (sky) {Skybox Points\\(optional)};
\draw[dasharrow] (sky) -- (params.east)
  node[midway, above, font=\footnotesize] {concat};

% ---- Annotation ----
\node[font=\footnotesize, left=0.3cm of params, align=right]
  {$\mathbf{x}$: position\\$\mathbf{f}$: SH colour\\$s$: log-scale\\$\mathbf{q}$: quaternion\\$o$: logit-opacity};

\end{tikzpicture}
\caption{Gaussian initialization from the fused point cloud. Each 3D point
$\mathbf{p}_i$ with colour $\mathbf{c}_i$ and camera-frame depth $d_i$ is
converted to a Gaussian primitive with position, SH colour coefficients,
log-scale (derived from depth and focal length), identity rotation, and
initial opacity. Optional skybox points are appended for open-environment
scenes.}
\label{fig:gaussian_init}
\end{figure*}


For each 3D point
$\mathbf{p}_i$ with colour $\mathbf{c}_i$ and camera-frame depth $d_i$,
a Gaussian primitive is created with the following parameters:

\begin{itemize}
    \item \textbf{Position} $\mathbf{x}_i = \mathbf{p}_i$ (world frame).
    \item \textbf{Colour (SH)}: The RGB colour is converted to zeroth-order
    spherical harmonics:
    $\mathbf{f}_i^{\text{dc}} = (\mathbf{c}_i - 0.5) / C_0$,
    where $C_0 = 0.28209$.
    \item \textbf{Scale}: $s_i = \log(\alpha \cdot d_i / f)$, where
    $\alpha$ is a scaling factor and $f$ is the average focal length.
    The log-parameterization ensures positivity after exponentiation.
    \item \textbf{Rotation}: Identity quaternion $[1, 0, 0, 0]$.
    \item \textbf{Opacity}: $\sigma^{-1}(0.1)$, the inverse sigmoid of 0.1.
\end{itemize}

For outdoor scenes, optional skybox points are added: $N_{\text{sky}}$
points uniformly distributed on a sphere of radius $10 \cdot R_{\text{sky}}$,
with a light-blue colour and higher initial opacity (0.7).

\subsection{Frame Addition and Depth Completion}

Each incoming aligned frame is added to the dataset. The RGB image is
converted from BGR (0--255) to RGB float (0--1). The fused depth map is
optionally completed by the SPNet TensorRT engine, which takes the RGB
image and sparse depth as input and produces a dense depth map.

After depth completion, the system selects patches where the completed
depth disagrees with the known LiDAR depth by less than 0.1\,m (mean
difference). In non-edge, positive-depth regions, additional 3D points
are back-projected from the completed depth at patch-level granularity,
enriching the point cloud with denser colour observations.

Keyframes are selected every $k$-th frame (default $k=5$). Keyframes are
added to the training set; non-keyframes to the test set for novel-view
evaluation.

\subsection{Map Extension}

When a new keyframe arrives and the map is already initialized, the
\texttt{extend} function adds new Gaussian primitives from the current
frame's point cloud. The process:

\begin{enumerate}
    \item Render the current viewpoint to obtain an alpha map
    (existing coverage).
    \item Project the new LiDAR points into the image plane and perform
    z-buffering to keep the nearest point per pixel.
    \item Filter out points that are: (a) outside the image bounds,
    (b) at non-positive depth, or (c) in already-opaque regions
    ($\alpha > 0.99$).
    \item Create Gaussian primitives for the surviving points with SH
    colour, log-scale from depth, identity rotation, and low initial
    opacity.
    \item Append to the existing Gaussian parameters via
    \texttt{densificationPostfix}, which also extends the optimizer state
    (Adam moments) for the new parameters.
\end{enumerate}

A memory safety cap (\texttt{max\_gaussians}) limits the total Gaussian
count on memory-constrained GPUs. When the budget is exceeded, new points
are randomly subsampled to fit the remaining budget.

\subsection{Map Optimization}

After each keyframe, the map is optimized by rendering a random subset of
training viewpoints and minimizing a photometric loss:

\begin{equation}
    \mathcal{L} = (1 - \lambda_{\text{dssim}}) \, \mathcal{L}_1
                + \lambda_{\text{dssim}} \, (1 - \text{SSIM})
                + \lambda_{\text{depth}} \, \mathcal{L}_{\text{depth}}
\end{equation}

where $\mathcal{L}_1$ is the L1 error between rendered and ground-truth
images, SSIM is computed via a fused CUDA kernel for efficiency, and
$\mathcal{L}_{\text{depth}}$ is the L1 error on depth where both rendered
and ground-truth depth are valid.

The optimizer is a custom \texttt{SparseGaussianAdam} that updates only
the Gaussians visible in the current rendered viewpoint, reducing
computational cost. The visibility mask is obtained from the rasterizer's
forward pass.

\subsection{Evaluation and Export}

After all frames are processed, the system evaluates visual quality on
both training and novel (test) viewpoints, computing PSNR, SSIM, and
LPIPS (using a pre-computed AlexNet-based LPIPS model loaded via TorchScript).

The final Gaussian map is exported as a binary PLY file with per-vertex
properties: position ($x, y, z$), SH DC coefficients, SH rest coefficients,
opacity, scale, and rotation. This format is compatible with standard
3DGS viewers such as the SIBR\_gaussianViewer and web-based renderers.
